\documentclass[12pt,a4paper]{article}
\usepackage[utf8]{inputenc} % inutile avec XeLaTeX/LuaLaTeX
\usepackage[T1]{fontenc}
\usepackage{amsmath,amssymb,mathrsfs,tikz,times,pifont}
\usepackage{enumitem}
\usepackage{multicol}
\usepackage{lmodern}
\newcommand\circitem[1]{%
\tikz[baseline=(char.base)]{
\node[circle,draw=gray, fill=red!55,
minimum size=1.2em,inner sep=0] (char) {#1};}}
\newcommand\boxitem[1]{%
\tikz[baseline=(char.base)]{
\node[fill=cyan,
minimum size=1.2em,inner sep=0] (char) {#1};}}
\setlist[enumerate,1]{label=\protect\circitem{\arabic*}}
\setlist[enumerate,2]{label=\protect\boxitem{\alph*}}
\everymath{\displaystyle}
\usepackage[left=1cm,right=1cm,top=1cm,bottom=1.7cm]{geometry}
\usepackage[colorlinks=true, linkcolor=blue, urlcolor=blue, citecolor=blue]{hyperref}
\usepackage{array,multirow}
\usepackage[most]{tcolorbox}
\usepackage{varwidth}
\usepackage{float}
\tcbuselibrary{skins,hooks}
\usetikzlibrary{patterns}

\newtcolorbox{exa}[2][]{enhanced,breakable,before skip=2mm,after skip=5mm,
colback=yellow!20!white,colframe=black!20!blue,boxrule=0.5mm,
attach boxed title to top left ={xshift=0.6cm,yshift*=1mm-\tcboxedtitleheight},
fonttitle=\bfseries,
title={#2},#1,
boxed title style={frame code={
\path[fill=tcbcolback!30!black]
([yshift=-1mm,xshift=-1mm]frame.north west)
arc[start angle=0,end angle=180,radius=1mm]
([yshift=-1mm,xshift=1mm]frame.north east)
arc[start angle=180,end angle=0,radius=1mm];
\path[left color=tcbcolback!60!black,right color = tcbcolback!60!black,
middle color = tcbcolback!80!black]
([xshift=-2mm]frame.north west) -- ([xshift=2mm]frame.north east)
[rounded corners=1mm]-- ([xshift=1mm,yshift=-1mm]frame.north east)
-- (frame.south east) -- (frame.south west)
-- ([xshift=-1mm,yshift=-1mm]frame.north west)
[sharp corners]-- cycle;
},interior engine=empty,
},interior style={top color=yellow!5}}

\usepackage{fancyhdr}
\usepackage{eso-pic}
\usepackage{tkz-tab}
\AddToShipoutPicture{
    \AtTextCenter{%
        \makebox[0pt]{\rotatebox{80}{\textcolor[gray]{0.7}{\fontsize{5cm}{5cm}\selectfont PGB}}}
    }
}

\usepackage{verbatim}
\usepackage{istgame}
\usepackage{color,soul}

\usepackage{amsmath}
\usepackage{amsfonts}
\usepackage{amssymb}
\usepackage{systeme}
\usepackage{tkz-tab}
\usepackage{tikz}
\usetikzlibrary{arrows}
\newcounter{exemple} % Compteur pour les questions

% Définir la commande pour afficher une question numérotée
\newcommand{\exemple}{%
  \refstepcounter{exemple}%
  \textbf{\textcolor{green}{Exemple \theexemple :}} \ignorespaces
}
%---------------------------------------
\definecolor{myorange}{rgb}{1.0, 0.8, 0.0}

% Définir un compteur pour les exercices d'application
\newcounter{exerciceapp}

% Définir la commande pour afficher un exercice d'application numéroté
\newcommand{\exerciceapp}{%
  \refstepcounter{exerciceapp}%
  \textbf{\textcolor{myorange}{Exercice d'application \theexerciceapp :}} \ignorespaces
}
%--------------------------------------
% Définir un compteur pour les exercices d'application
\newcounter{correction}

% Définir la commande pour afficher un correction exercice d'application numéroté
\newcommand{\correction}{%
  \refstepcounter{correction}%
  \textbf{\textcolor{myorange}{Correction \thecorrection :}} \ignorespaces
}
%--------------------------------------
% Commande pour ajouter du texte en arrière-plan
\usepackage{fancyhdr}
\usepackage{eso-pic}
\usepackage{diagbox}
\usepackage{xcolor} % Pour ajouter les couleurs

% Définition d'un rouge foncé similaire à l'encre du cahier
\definecolor{darkred}{rgb}{0.7, 0, 0}
\definecolor{darkgreen}{rgb}{0, 0.7, 0}
%\usepackage{tkz-tab}
\AddToShipoutPicture{
    \AtTextCenter{%
        \makebox[0pt]{\rotatebox{80}{\textcolor[gray]{0.7}{\fontsize{5cm}{5cm}\selectfont PGB}}}
    }
}
%This command takes a colour as an optional argument; the default colour is black.
\usetikzlibrary{shapes.geometric,fit}
\newcommand{\myul}[2][black]{\setulcolor{#1}\ul{#2}\setulcolor{black}}
\newcommand\tab[1][1cm]{\hspace*{#1}}

\begin{document}
% En-tête personnalisée
\begin{center}
    \Large\textbf{\underline{\textcolor{red}{Limites-Continuités}}}\\[-0.1cm]
    \normalsize\textbf{Prof : M. BA} \hfill \textbf{Classe : 1er S2}\\[-0.1cm]
    \textbf{Année scolaire : 2026 -- 2027}
\end{center}

\section*{\underline{\textcolor{darkred}{INTRODUCTION}}}

\begin{tcolorbox}[
    colback=white,
    colframe=red,
    boxrule=0.5pt,
    sharp corners,
    left=8pt,
    right=8pt,
    top=8pt,
    bottom=8pt
]

La courbe ci-dessous représente une fonction \(f:x\mapsto f(x)\).

\begin{center}
\begin{tikzpicture}[scale=1.1]

    % Axes
    \draw[->] (-0.2,0) -- (7,0);
    \draw[->] (0,-0.3) -- (0,4.5);

    % Graduations sur l'axe des abscisses
    \foreach \x in {1,2,3,4,5,6}
        \draw (\x,0.05) -- (\x,-0.05);

    % Graduations sur l'axe des ordonnées
    \foreach \y in {1,2,3,4}
        \draw (0.05,\y) -- (-0.05,\y);

    % Points I et J
    \fill (0,1) circle (1.5pt);
    \node[left] at (0,1) {\(I\)};

    \fill (1,0) circle (1.5pt);
    \node[below] at (1,0) {\(J\)};

    % Origine
    \node[below left] at (0,0) {\(O\)};

    % Courbe
    \draw[blue!60]
        plot[
            domain=0.18:6.5,
            samples=250
        ]
        (\x,{1/\x});

\end{tikzpicture}
\end{center}

\medskip

\begin{flushleft}
\textit{
Considérons la fonction \(f\) représentée ci-dessus et posons-nous les questions suivantes :
}
\end{flushleft}

\begin{itemize}
    \item \textit{Que devient \(f(x)\) lorsque la variable \(x\) s'éloigne vers \(+\infty\) ?}
    
    \item \textit{Que devient \(f(x)\) lorsque la variable \(x\) se rapproche de \(0\) ?}
\end{itemize}

\textit{
C'est en essayant de répondre à ce type de question que la notion de limite a été définie.
}

\medskip

\textit{
Ainsi, si on observe la courbe de \(f\), on voit que lorsque \(x\) s'éloigne vers \(+\infty\), \(f(x)\)
se rapproche de \(0\) ;
}

\textit{
on dit que la limite lorsque \(x\) tend vers \(+\infty\) de \(f\) est égale à \(0\).
}

\textit{
De même, si on observe la courbe, on voit que lorsque \(x\) se rapproche de \(0\) tout en restant
à droite de \(0\), \(f(x)\) tend vers \(+\infty\) ;
}

\textit{
on dit que la limite à droite de \(0\) de \(f\) est égale à \(+\infty\).
}

\end{tcolorbox}

\end{document}